\documentclass[../../main.tex]{subfiles}
\begin{document}

\begin{center}
  {\fontsize{48pt}{58pt}\selectfont \textbf{1964}\normalsize\textbf{年度}}
\end{center}

\begin{table}[htb]
  \centering
  \caption{出題テーマと難易度}
  \begin{tabular}{|c|p{9cm}|c|}
    \hline
    \textbf{問題番号} & \textbf{テーマ} & \textbf{難易度} \\
    \hline
    第1問 & 相加平均・相乗平均の不等式と順列 & \\
    \hline
    第2問 & 多項式の最大公約数と最小公倍数 & \\
    \hline
    第3問 & 三角形内の3線分の長さと三角形の成立条件 & \\
    \hline
    第4問 & 放物線の2接線が一定の角をなす点の軌跡 & \\
    \hline
    第5問 & 三角級数の収束と積分の交換 & \\
    \hline
    第6問 & 余弦曲線内の図形の最大面積 & \\
    \hline
  \end{tabular}
\end{table}

\begin{table}[htb]
  \centering
  \caption{解答}
  \renewcommand{\arraystretch}{1.4}
  \begin{tabular}{|c|c|p{9cm}|}
    \hline
    \textbf{問題番号} & \textbf{小問} & \textbf{答え} \\
    \hline
    第1問 & (1) & 不等式は成立する（証明問題） \\
    \cline{2-3}
    & (2) & 和は$4$以上（証明問題） \\
    \hline
    第2問 & -- & $A=x-1,B=x+6,C=x-6$とすると，$(x+3)$を除いた組は$\{A^2,AB,BC\}$または$\{A^2,AC,BC\}$ \\
    \hline
    第3問 & (1) & $RP\le QR\le BQ$ \\
    \cline{2-3}
    & (2) & $AC$の中点を$D$とすると，$P$は線分$DC$上（$D$を除く） \\
    \hline
    第4問 & -- & $x^2-3\left(y+\dfrac{5}{12}\right)^2=\dfrac{1}{9}$ \\
    \hline
    第5問 & (イ) & $\sin a_n$と$\cos a_n$がともに収束すること \\
    \cline{2-3}
    & (ロ) & $\alpha=\lim\sin a_n,\beta=\lim\cos a_n$とすると$F(x)=\dfrac{\beta\sin x+\alpha\cos x-\sin(x+a_1)}{2}$．二つの積分値は等しい \\
    \hline
    第6問 & -- & $A<B=C$ \\
    \hline
  \end{tabular}
\end{table}

\paragraph{試験形態}
6問．試験時間は要確認．

\end{document}
